\section{System Architecture}
\label{ch:architecture}

\subsection{Architecture and data flow}

The application is organised around six responsibilities: TORCS simulation, a Python runner, agent implementations, telemetry and storage, the desktop GUI, and separate training/evaluation scripts. \Cref{fig:architecture-overview} depicts the agent-environment loop and evidence flow. TORCS provides observations; the selected agent produces steering, throttle, brake and gear decisions; the runner applies them and records the resulting state; storage persists both run-level outcomes and time-series telemetry; and the GUI reloads those records for live and post-race inspection.

\begin{figure}[htbp]
  \centering
  \includegraphics[width=\textwidth]{architecture-overview.png}
  \caption{System architecture of Enhanced AI Racing. The novice-facing application launches and explains fixed policies, while offline training and formal evaluation remain separate from the runtime control loop and preserve auditable evidence.}
  \label{fig:architecture-overview}
\end{figure}

This separation is not merely stylistic. Keeping the runner independent of agent logic means that a new controller can be compared without rebuilding the interface or database. Keeping the GUI independent of training means a novice can run a protected policy without accidentally modifying a checkpoint. Keeping formal evaluation scripts separate from GUI demonstrations protects evidence provenance: a visible application run is useful for explanation, whereas a fixed-policy repeated evaluation is stronger evidence of reliability.

The storage layer has two granularities. Run-level records contain agent identity, track, start time, completion status, lap timing and termination information. Telemetry records contain the values needed to reconstruct how that outcome occurred, including speed, track position and controls. The distinction supports efficient use of the interface: a comparison table can load summary records quickly, while a review screen retrieves the detailed trace only when a user asks why a run failed. It also aligns the data model with the dissertation's evidence hierarchy, where every aggregate claim should remain traceable to individual trials.

\subsection{Deployment and the novice workflow}

The release package contains the executable, its required internal runtime files, TORCS, final Agent~7 and Agent~8 policies, Dyna-Q policies, racing lines and compact results evidence. The user extracts the complete Enhanced AI Racing folder and opens \texttt{EnhancedAIRacing.exe}. The dashboard starts without a terminal; TORCS starts only after the user selects an agent and presses Start. A package smoke test verifies that the dashboard opens, expected agents appear, final policies load, TORCS launches and a completed run is visible in history.

This design directly responds to the brief's request for a launcher and beginner-facing wrapper. It also creates a clear pedagogical path: users can observe a weak random baseline, compare it with explicit rules, inspect a racing line and then encounter learned policies. The design is accessibility-oriented, including readable labels, results summaries, language support and colour choices that do not make colour the sole carrier of meaning \parencite{wcag22}. Its limitation is that the dissertation evaluates functional accessibility and explanatory design, not measured learning gains from a controlled user study.

The architecture deliberately avoids turning the novice application into a training console. Long-running learning scripts remain separate because they have different failure modes, time requirements and vocabulary from the task of observing a finished policy. The GUI exposes enough evidence to understand the result, while advanced details such as policy path and model contract are available in the research view. This division reduces cognitive load without preventing technical audit, and it avoids the misleading impression that pressing Start is equivalent to training a reinforcement-learning agent from scratch.

